\documentclass[11pt]{article}
\usepackage[a4paper,margin=2.0cm]{geometry}
\usepackage{amsmath,amssymb,booktabs,graphicx}
\usepackage[T1]{fontenc}
\usepackage{hyperref}
% generated by qbscreen/make_numbers_v5.py — do not edit
\newcommand{\Dcontact}{-843}
\newcommand{\DeltaCorrBenzyl}{2.3}
\newcommand{\DeltaCorrMethyl}{2.1}
\newcommand{\DeltaDFTmax}{2.99}
\newcommand{\DeltaDFTmin}{2.47}
\newcommand{\DeltaGap}{0.75}
\newcommand{\DeltaHundredthOne}{0.53}
\newcommand{\DeltaHundredthOneEnv}{0.53}
\newcommand{\DeltaHundredthOneNoD}{0.53}
\newcommand{\DeltaHundredthThousand}{0.34}
\newcommand{\DeltaHundredthThousandEnv}{0.34}
\newcommand{\DeltaHundredthThousandNoD}{0.34}
\newcommand{\DeltaLit}{1.61}
\newcommand{\DeltaMaxHi}{0.86}
\newcommand{\DeltaMaxLo}{0.62}
\newcommand{\DeltaTenthOne}{0.44}
\newcommand{\DeltaTenthOneEnv}{0.47}
\newcommand{\DeltaTenthOneNoD}{0.47}
\newcommand{\DeltaTenthThousand}{0.25}
\newcommand{\DeltaTenthThousandEnv}{0.28}
\newcommand{\DeltaTenthThousandNoD}{0.28}
\newcommand{\EamExp}{1.30}
\newcommand{\Econtact}{-26}
\newcommand{\EflExp}{-0.313}
\newcommand{\EflNeeded}{0.44}
\newcommand{\ErrAm}{-0.57}
\newcommand{\ErrFl}{-0.84}
\newcommand{\ErrNet}{0.27}
\newcommand{\FkPhundred}{0.02}
\newcommand{\FkPhundredNoD}{0.013}
\newcommand{\FkPlow}{1.4}
\newcommand{\FkPlowNoD}{14}
\newcommand{\FkPtop}{1.1\times10^{-7}}
\newcommand{\FkPtopNoD}{1.1\times10^{-7}}
\newcommand{\Fmax}{1.4}
\newcommand{\FmaxNoD}{14}
\newcommand{\GatDmaxOne}{1.1\times10^{-7}}
\newcommand{\GatDmaxOneEnv}{1.1\times10^{-7}}
\newcommand{\GatDmaxOneNoD}{1.1\times10^{-7}}
\newcommand{\GatDmaxThousand}{1.1\times10^{-7}}
\newcommand{\GatDmaxThousandEnv}{1.1\times10^{-7}}
\newcommand{\GatDmaxThousandNoD}{1.1\times10^{-7}}
\newcommand{\GzeroOne}{1.4}
\newcommand{\GzeroOneEnv}{14}
\newcommand{\GzeroOneNoD}{14}
\newcommand{\GzeroThousand}{1.4}
\newcommand{\GzeroThousandEnv}{14}
\newcommand{\GzeroThousandNoD}{14}
\newcommand{\aBenzOne}{257}
\newcommand{\aBenzTwo}{196}
\newcommand{\aTrpOne}{60}
\newcommand{\aTrpTwo}{56}
\newcommand{\cryDcoup}{-11.4}
\newcommand{\cryMax}{0.63}
\newcommand{\cryMidMax}{0.63}
\newcommand{\cryMidMin}{0.17}
\newcommand{\cryMin}{0.0015}
\newcommand{\cryTtwoRatio}{0.6}
\newcommand{\kTln}{0.037}
\newcommand{\kfLit}{6\times10^{-14}}
\newcommand{\nCells}{195}
\newcommand{\nCellsNoD}{195}
\newcommand{\nContactExceeds}{74}
\newcommand{\nEdge}{0}
\newcommand{\nEdgeNoD}{0}
\newcommand{\nScan}{112}
\newcommand{\nScanValid}{69}
\newcommand{\ratioContactExceeds}{8}
\newcommand{\tContactMax}{106}
\newcommand{\tContactMed}{53}

\renewcommand{\thetable}{S\arabic{table}}
\renewcommand{\thesection}{S\arabic{section}}
\renewcommand{\theequation}{S\arabic{equation}}

\title{Supporting Information:\\ Catalytic competence bounds the magnetic field
effect of a thermally formed radical pair}
\author{Hikaru Wakaura and Taiki Tanimae}
\date{}

\begin{document}
\maketitle

\section{Software and environment}
Spin dynamics: \texttt{qbscreen} (this work) with NumPy and SciPy; coherent
yields from the Sylvester equation
$H_{\rm eff}X - XH_{\rm eff}^\dagger = -i\rho_0$ with
$H_{\rm eff} = 2\pi H - \tfrac{i}{2}(k_SP_S + k_P)$ and $X = \int_0^\infty\rho\,dt$;
yields with electron spin relaxation from the Liouville-space generator. The two
agree to $10^{-8}$ in the absence of relaxation, and every yield pair sums to
one. Electronic structure: PySCF~2.12.1 (B3LYP; def2-SVP for couplings and
energies, def2-TZVP for hyperfine couplings; RI-J; PCM), geometries from
GFN2-xTB~6.7.1. Every table below is generated from the stored results by
\texttt{qbscreen/make\_si\_v5.py}.

\section{Proof of the bound $\Phi_P\le4k_P/k_S$}
Write the generator of eq~2 of the main text as
$\mathcal L = \mathcal L_0 - k_P$, where $\mathcal L_0\rho = -2\pi i[H,\rho]
-\tfrac{k_S}{2}\{P_S,\rho\} + \mathcal D\rho$ and $\mathcal D$ is any unital
relaxation superoperator (for single-electron dephasing
$\mathcal D\rho = \sum_i (2/T_2^e)(S_{iz}\rho S_{iz}-\rho/4)$). Under a constant
source $s\,P_S/N$ into the singlet ($N = {\rm Tr}P_S$), the steady state is
\begin{equation}
\rho_{\rm ss}(k_P) = \int_0^\infty e^{-k_Pt}\,e^{\mathcal L_0t}\,(sP_S/N)\,dt ,
\end{equation}
because $k_P$ commutes with $\mathcal L_0$. The trace of $e^{\mathcal L_0t}(P_S/N)$
is non-negative and non-increasing, so
${\rm Tr}\rho_{\rm ss}(k_P)\le{\rm Tr}\rho_{\rm ss}(0)$. At $k_P=0$ let $\Pi$
be the projector onto the smallest subspace that contains the range of $P_S$ and
is invariant under $H$ and $\mathcal D$. Then $c\,\Pi$ is stationary under the
coherent and relaxation terms (the steady state on $\Pi$ is unique, because by
minimality $\Pi$ contains no invariant subspace orthogonal to $P_S$), and the
balance of source and singlet decay,
$sP_S/N = c\,k_SP_S$, fixes $c = s/(k_SN)$; hence
${\rm Tr}\rho_{\rm ss}(0) = s\,{\rm dim}\,\Pi/(k_SN)\le 4s/k_S$, since
${\rm dim}\,\Pi\le 4N$. With ${\rm Tr}\rho_{\rm ss}(k_P) = s\Phi_P/k_P$ this is
$\Phi_P\le4k_P/k_S$. With $s = k_f[{\rm E\cdot S}]$ and
$k_f = k_Se^{-\Delta_S/k_BT}$ the flux through the pair is at most
$k_P\,e^{-\Delta/k_BT}[{\rm E\cdot S}]$. The inequality was tested on 300 random spin systems (0--3 nuclei of spin
$\tfrac12$ or 1 with couplings of 1--300~MHz, exchange $10^{-1}$--$10^4$~MHz of
either sign, random axial dipolar tensors, fields and rates, and in part with
$T_2^e$ from 1~ns to 10~$\mu$s); the largest value of $\Phi_Pk_S/4k_P$ found is
0.99995, approached when singlet--triplet mixing is complete. The self-check in
\texttt{competence\_sensitivity} repeats the test on 40 systems.

\section{Worst-case map: numerical details}
For each $(k_S,k_P)$ the exchange coupling is scanned on 119 points with
$\sinh$ spacing (40~MHz scale), symmetric about zero and extending to
$\max(4000~{\rm MHz},5(k_S+k_P)/2\pi)$; the three largest maxima are refined
three times on 21-point subgrids. The field direction runs over the pole and
four azimuths at each of four polar angles on a hemisphere (17 directions; the
response is inversion symmetric). A cell whose maximum lies within 5\,\% of the
range limit is recomputed with a range sixteen times the position of the
maximum. Tables~\ref{tab:mapc}--\ref{tab:maps} list $F$ in percent. The bound $G$ of the main text uses the cell-by-cell maximum over all three tables.

\begin{table}[h]
\caption{$F(k_S,k_P)$ (\%) with the contact dipolar tensor. Rates in
$\mu$s$^{-1}$.}
\label{tab:mapc}
\centering
\resizebox{\textwidth}{!}{\scriptsize\setlength{\tabcolsep}{2pt}\begin{tabular}{rrrrrrrrrrrrrrrr}
\toprule
$\log k_P \backslash \log k_S$ & -2.0 & -1.5 & -1.0 & -0.5 & 0.0 & 0.5 & 1.0 & 1.5 & 2.0 & 2.5 & 3.0 & 3.5 & 4.0 & 4.5 & 5.0 \\
\midrule
-2.0 & 0.76 & 1.4 & 1.3 & 0.71 & 1.2 & 1.1 & 0.81 & 0.49 & 0.33 & 0.18 & 0.14 & 0.14 & 0.14 & 0.14 & 0.14 \\
-1.5 & 0.3 & 0.68 & 1.1 & 0.95 & 0.5 & 0.96 & 0.93 & 0.67 & 0.49 & 0.32 & 0.18 & 0.14 & 0.14 & 0.14 & 0.14 \\
-1.0 & 0.073 & 0.22 & 0.49 & 0.75 & 0.64 & 0.45 & 0.82 & 0.66 & 0.66 & 0.48 & 0.3 & 0.16 & 0.14 & 0.14 & 0.14 \\
-0.5 & 0.022 & 0.067 & 0.19 & 0.42 & 0.64 & 0.6 & 0.36 & 0.38 & 0.65 & 0.64 & 0.41 & 0.21 & 0.14 & 0.14 & 0.12 \\
0.0 & 0.0065 & 0.02 & 0.061 & 0.17 & 0.4 & 0.6 & 0.53 & 0.23 & 0.36 & 0.58 & 0.44 & 0.17 & 0.13 & 0.12 & 0.059 \\
0.5 & 0.0016 & 0.0051 & 0.016 & 0.048 & 0.14 & 0.3 & 0.39 & 0.2 & 0.13 & 0.29 & 0.27 & 0.11 & 0.11 & 0.057 & 0.042 \\
1.0 & 2.0\text{e-}4 & 6.4\text{e-}4 & 0.002 & 0.0061 & 0.018 & 0.042 & 0.061 & 0.052 & 0.03 & 0.058 & 0.064 & 0.05 & 0.048 & 0.043 & 0.035 \\
1.5 & 1.2\text{e-}5 & 3.7\text{e-}5 & 1.2\text{e-}4 & 3.6\text{e-}4 & 0.0011 & 0.0031 & 0.008 & 0.017 & 0.02 & 0.016 & 0.031 & 0.044 & 0.044 & 0.034 & 0.017 \\
2.0 & 8.7\text{e-}7 & 2.8\text{e-}6 & 8.7\text{e-}6 & 2.7\text{e-}5 & 8.6\text{e-}5 & 2.7\text{e-}4 & 7.9\text{e-}4 & 0.002 & 0.004 & 0.0058 & 0.011 & 0.016 & 0.02 & 0.013 & 0.0038 \\
2.5 & 4.3\text{e-}8 & 1.4\text{e-}7 & 4.3\text{e-}7 & 1.4\text{e-}6 & 4.3\text{e-}6 & 1.4\text{e-}5 & 4.2\text{e-}5 & 1.3\text{e-}4 & 3.7\text{e-}4 & 9.5\text{e-}4 & 0.0021 & 0.0037 & 0.0033 & 0.0013 & 2.4\text{e-}4 \\
3.0 & 4.3\text{e-}9 & 1.4\text{e-}8 & 4.3\text{e-}8 & 1.4\text{e-}7 & 4.3\text{e-}7 & 1.4\text{e-}6 & 4.3\text{e-}6 & 1.3\text{e-}5 & 4.0\text{e-}5 & 1.1\text{e-}4 & 2.2\text{e-}4 & 2.6\text{e-}4 & 1.4\text{e-}4 & 3.2\text{e-}5 & 4.5\text{e-}6 \\
3.5 & 9.9\text{e-}11 & 3.1\text{e-}10 & 9.9\text{e-}10 & 3.1\text{e-}9 & 9.9\text{e-}9 & 3.1\text{e-}8 & 9.9\text{e-}8 & 3.1\text{e-}7 & 9.3\text{e-}7 & 2.5\text{e-}6 & 5.2\text{e-}6 & 5.2\text{e-}6 & 2.1\text{e-}6 & 2.4\text{e-}7 & 3.3\text{e-}8 \\
4.0 & 8.9\text{e-}13 & 2.5\text{e-}12 & 7.4\text{e-}12 & 2.3\text{e-}11 & 7.3\text{e-}11 & 2.3\text{e-}10 & 7.2\text{e-}10 & 2.3\text{e-}9 & 7.1\text{e-}9 & 2.1\text{e-}8 & 5.7\text{e-}8 & 1.1\text{e-}7 & 1.1\text{e-}7 & 3.2\text{e-}8 & 2.9\text{e-}9 \\
\bottomrule
\end{tabular}
}
\end{table}

\begin{table}[h]
\caption{$F(k_S,k_P)$ (\%) without dipolar coupling.}
\label{tab:mapn}
\centering
\resizebox{\textwidth}{!}{\scriptsize\setlength{\tabcolsep}{2pt}\begin{tabular}{rrrrrrrrrrrrrrrr}
\toprule
$\log k_P \backslash \log k_S$ & -2.0 & -1.5 & -1.0 & -0.5 & 0.0 & 0.5 & 1.0 & 1.5 & 2.0 & 2.5 & 3.0 & 3.5 & 4.0 & 4.5 & 5.0 \\
\midrule
-2.0 & 6.3 & 12 & 14 & 13 & 10 & 7.5 & 6 & 5.2 & 4.6 & 3.7 & 2.3 & 1.1 & 0.4 & 0.13 & 0.05 \\
-1.5 & 2.5 & 6.2 & 11 & 14 & 13 & 9.6 & 7.1 & 5.3 & 3.8 & 2.3 & 1.1 & 0.4 & 0.13 & 0.05 & 0.04 \\
-1.0 & 0.85 & 2.4 & 6 & 11 & 13 & 12 & 8.4 & 5.4 & 2.7 & 1.1 & 0.4 & 0.13 & 0.05 & 0.04 & 0.039 \\
-0.5 & 0.26 & 0.8 & 2.3 & 5.5 & 9.7 & 11 & 9 & 5.2 & 2 & 0.54 & 0.15 & 0.051 & 0.043 & 0.039 & 0.03 \\
0.0 & 0.069 & 0.22 & 0.66 & 1.9 & 4.4 & 7.2 & 7.4 & 4.4 & 1.6 & 0.42 & 0.092 & 0.051 & 0.041 & 0.03 & 0.019 \\
0.5 & 0.014 & 0.044 & 0.14 & 0.41 & 1.1 & 2.5 & 3.6 & 2.7 & 1.1 & 0.37 & 0.1 & 0.047 & 0.03 & 0.018 & 0.011 \\
1.0 & 0.0014 & 0.0044 & 0.014 & 0.043 & 0.13 & 0.34 & 0.66 & 0.75 & 0.48 & 0.2 & 0.069 & 0.03 & 0.018 & 0.011 & 0.0057 \\
1.5 & 5.8\text{e-}5 & 1.8\text{e-}4 & 5.8\text{e-}4 & 0.0018 & 0.0056 & 0.017 & 0.043 & 0.083 & 0.096 & 0.062 & 0.029 & 0.017 & 0.01 & 0.0055 & 0.0021 \\
2.0 & 1.7\text{e-}6 & 5.4\text{e-}6 & 1.7\text{e-}5 & 5.4\text{e-}5 & 1.7\text{e-}4 & 5.2\text{e-}4 & 0.0015 & 0.004 & 0.0077 & 0.012 & 0.013 & 0.0091 & 0.005 & 0.0019 & 4.9\text{e-}4 \\
2.5 & 1.2\text{e-}7 & 3.8\text{e-}7 & 1.2\text{e-}6 & 3.8\text{e-}6 & 1.2\text{e-}5 & 3.7\text{e-}5 & 1.2\text{e-}4 & 3.5\text{e-}4 & 9.6\text{e-}4 & 0.0021 & 0.0033 & 0.0028 & 0.0014 & 3.9\text{e-}4 & 6.5\text{e-}5 \\
3.0 & 4.8\text{e-}9 & 1.5\text{e-}8 & 4.8\text{e-}8 & 1.5\text{e-}7 & 4.8\text{e-}7 & 1.5\text{e-}6 & 4.8\text{e-}6 & 1.5\text{e-}5 & 4.4\text{e-}5 & 1.1\text{e-}4 & 2.2\text{e-}4 & 2.5\text{e-}4 & 1.2\text{e-}4 & 2.7\text{e-}5 & 3.7\text{e-}6 \\
3.5 & 5.8\text{e-}11 & 1.8\text{e-}10 & 5.8\text{e-}10 & 1.8\text{e-}9 & 5.8\text{e-}9 & 1.8\text{e-}8 & 5.8\text{e-}8 & 1.8\text{e-}7 & 5.4\text{e-}7 & 1.5\text{e-}6 & 3.0\text{e-}6 & 3.2\text{e-}6 & 1.4\text{e-}6 & 3.0\text{e-}7 & 5.7\text{e-}8 \\
4.0 & 1.0\text{e-}12 & 2.6\text{e-}12 & 7.3\text{e-}12 & 2.2\text{e-}11 & 6.9\text{e-}11 & 2.2\text{e-}10 & 6.9\text{e-}10 & 2.2\text{e-}9 & 6.7\text{e-}9 & 2.0\text{e-}8 & 5.4\text{e-}8 & 1.1\text{e-}7 & 1.0\text{e-}7 & 3.1\text{e-}8 & 2.7\text{e-}9 \\
\bottomrule
\end{tabular}
}
\end{table}

\begin{table}[h]
\caption{$F(k_S,k_P)$ (\%) with the contact dipolar tensor scaled by the factor
shown, on the onward rates that set the thresholds of $G$.}
\label{tab:maps}
\centering
\resizebox{\textwidth}{!}{\scriptsize\setlength{\tabcolsep}{2pt}\begin{tabular}{rrrrrrrrrrrrrrrrr}
\toprule
scale & $\log k_P \backslash \log k_S$ & -2.0 & -1.5 & -1.0 & -0.5 & 0.0 & 0.5 & 1.0 & 1.5 & 2.0 & 2.5 & 3.0 & 3.5 & 4.0 & 4.5 & 5.0 \\
\midrule
0.1 & 1.0 & 2.5\text{e-}4 & 7.8\text{e-}4 & 0.0024 & 0.0076 & 0.022 & 0.058 & 0.11 & 0.12 & 0.082 & 0.05 & 0.034 & 0.021 & 0.016 & 0.011 & 0.0065 \\
0.1 & 1.5 & 1.7\text{e-}5 & 5.3\text{e-}5 & 1.7\text{e-}4 & 5.2\text{e-}4 & 0.0016 & 0.0049 & 0.013 & 0.027 & 0.034 & 0.027 & 0.02 & 0.014 & 0.0094 & 0.0064 & 0.0027 \\
0.1 & 2.0 & 1.2\text{e-}6 & 3.7\text{e-}6 & 1.2\text{e-}5 & 3.7\text{e-}5 & 1.1\text{e-}4 & 3.5\text{e-}4 & 0.001 & 0.0025 & 0.0049 & 0.0084 & 0.01 & 0.0079 & 0.0057 & 0.0025 & 6.3\text{e-}4 \\
0.3 & 1.0 & 2.0\text{e-}4 & 6.2\text{e-}4 & 0.002 & 0.006 & 0.018 & 0.044 & 0.08 & 0.096 & 0.063 & 0.044 & 0.031 & 0.022 & 0.019 & 0.019 & 0.015 \\
0.3 & 1.5 & 1.8\text{e-}5 & 5.7\text{e-}5 & 1.8\text{e-}4 & 5.7\text{e-}4 & 0.0018 & 0.0053 & 0.014 & 0.028 & 0.03 & 0.022 & 0.026 & 0.019 & 0.016 & 0.014 & 0.0066 \\
0.3 & 2.0 & 1.1\text{e-}6 & 3.5\text{e-}6 & 1.1\text{e-}5 & 3.5\text{e-}5 & 1.1\text{e-}4 & 3.3\text{e-}4 & 9.5\text{e-}4 & 0.0022 & 0.0036 & 0.0056 & 0.0092 & 0.011 & 0.0096 & 0.0051 & 0.0014 \\
3 & 1.0 & 8.4\text{e-}5 & 2.7\text{e-}4 & 8.3\text{e-}4 & 0.0026 & 0.0075 & 0.018 & 0.032 & 0.033 & 0.037 & 0.033 & 0.042 & 0.071 & 0.083 & 0.076 & 0.045 \\
3 & 1.5 & 1.2\text{e-}5 & 3.8\text{e-}5 & 1.2\text{e-}4 & 3.7\text{e-}4 & 0.0011 & 0.0033 & 0.0077 & 0.013 & 0.019 & 0.031 & 0.032 & 0.035 & 0.046 & 0.037 & 0.016 \\
3 & 2.0 & 8.4\text{e-}7 & 2.7\text{e-}6 & 8.4\text{e-}6 & 2.7\text{e-}5 & 8.3\text{e-}5 & 2.6\text{e-}4 & 7.5\text{e-}4 & 0.0021 & 0.0058 & 0.012 & 0.018 & 0.02 & 0.022 & 0.013 & 0.0037 \\
\bottomrule
\end{tabular}
}
\end{table}

\section{Dipolar tensor}
The tensor is the point-dipole tensor summed over pairs of atoms weighted by
their Löwdin spin populations (UKS B3LYP/def2-SVP on each radical at its position
in the complex), for methylamine on the flavin si face with N5--N = 3.5~\AA\
(Table~\ref{tab:dip}). Summing over the spin distributions halves $D$ relative
to point dipoles at the spin centroids.

\begin{table}[h]
\caption{Dipolar coupling of the contact pair.}
\label{tab:dip}
\centering
\begin{tabular}{lrrl}
\toprule
model & $D$ (MHz) & $E$ (MHz) & principal values of $\mathsf T$ (MHz) \\
\midrule
spin-population tensor & -843 & -26 & -1124, 511, 613 \\
point dipole at spin centroids & -1642 & 0 & -- \\
point dipole at N5 and N & -1817 & 0 & -- \\
\bottomrule
\end{tabular}

\end{table}

\section{Hyperfine couplings}
Fermi-contact couplings (B3LYP/def2-TZVP at GFN2-xTB geometries), all nuclei
with $|a_{\rm iso}|\ge1$~MHz (Table~\ref{tab:hfc}). For the flavin the main text
uses the literature values for N5 and N10;\cite{Lee2014} the computed values
are within 16\,\% and 13\,\% of them. For 3-methylindole$^{\bullet+}$ the two
large methyl couplings stand in for the tryptophan $\beta$-protons; their
values depend on the side-chain conformation.

\begin{table}[h]
\caption{Isotropic hyperfine couplings.}
\label{tab:hfc}
\centering
\small\begin{tabular}{llrr}
\toprule
radical & nucleus (atom index) & $a_{\rm iso}$ (MHz) & $a_{\rm iso}$ ($\mu$T) \\
\midrule
lumiflavin$^{\bullet-}$ & H29 & +17.9 & +639 \\
 & H30 & +17.8 & +636 \\
 & N4 & +12.4 & +441 \\
 & H25 & +11.3 & +403 \\
 & H26 & +11.3 & +402 \\
 & H22 & -10.5 & -374 \\
 & H20 & -5.3 & -190 \\
 & H21 & -5.3 & -190 \\
 & N13 & +4.6 & +165 \\
 & H27 & +1.4 & +50 \\
 & N8 & -1.0 & -36 \\
methylaminium$^{\bullet+}$ & H2 & +363.1 & +12955 \\
 & H6 & +82.7 & +2951 \\
 & H5 & +81.9 & +2923 \\
 & H3 & -59.3 & -2116 \\
 & H4 & -59.3 & -2115 \\
 & N1 & +30.5 & +1090 \\
benzylaminium$^{\bullet+}$ & H10 & +256.9 & +9167 \\
 & H11 & +196.2 & +7001 \\
 & N0 & +24.8 & +885 \\
 & H8 & -23.0 & -821 \\
 & H9 & -20.7 & -737 \\
 & H14 & -12.3 & -439 \\
 & H16 & -7.8 & -278 \\
 & H13 & -3.2 & -115 \\
 & H15 & +2.4 & +85 \\
3-methylindole$^{\bullet+}$ & H12 & +60.1 & +2143 \\
 & H11 & +55.5 & +1980 \\
 & H18 & -14.3 & -511 \\
 & H13 & -13.1 & -467 \\
 & H14 & -13.1 & -466 \\
 & H16 & -9.8 & -351 \\
 & H15 & -4.9 & -175 \\
 & N3 & +4.8 & +172 \\
 & H17 & +2.7 & +95 \\
\bottomrule
\end{tabular}

\end{table}

\section{Electronic coupling}
Fragment-orbital couplings $t = \langle{\rm LUMO_{Fl}}|F|{\rm HOMO_{am}}\rangle$
(B3LYP/def2-SVP) for methylamine and benzylamine stacked on the flavin si face
(five orientations) or approaching its edge (three orientations), at N5--N
distances of 3.0--6.0~\AA. Geometries with any intermolecular contact below
2.0~\AA\ are excluded (Table~\ref{tab:t}). The edge approach of benzylamine
clashes at every distance below 6~\AA\ and is not used.

\begin{table}[h]
\caption{Coupling scan. ``valid'': non-clashing geometries of those computed.}
\label{tab:t}
\centering
\small\begin{tabular}{llrrrrr}
\toprule
amine & approach & $r_{\rm N5N}$ (\AA) & valid & median $|t|$ (meV) & max $|t|$ (meV) & min contact (\AA) \\
\midrule
methylamine & face & 3.0 & 5/5 & 184.6 & 311.6 & 2.32 \\
methylamine & face & 3.5 & 5/5 & 52.8 & 105.6 & 2.72 \\
methylamine & face & 4.0 & 5/5 & 24.6 & 37.9 & 3.12 \\
methylamine & face & 4.5 & 5/5 & 5.1 & 31.1 & 3.53 \\
methylamine & face & 5.0 & 5/5 & 1.9 & 14.7 & 3.96 \\
methylamine & face & 5.5 & 5/5 & 1.0 & 4.2 & 4.40 \\
methylamine & face & 6.0 & 5/5 & 0.3 & 1.1 & 4.85 \\
methylamine & edge & 3.0 & 0/3 & -- & -- & 0.77 \\
methylamine & edge & 3.5 & 0/3 & -- & -- & 0.93 \\
methylamine & edge & 4.0 & 0/3 & -- & -- & 1.29 \\
methylamine & edge & 4.5 & 2/3 & 97.0 & 170.8 & 1.71 \\
methylamine & edge & 5.0 & 3/3 & 23.1 & 81.0 & 2.17 \\
methylamine & edge & 5.5 & 3/3 & 8.6 & 35.0 & 2.65 \\
methylamine & edge & 6.0 & 3/3 & 3.2 & 12.7 & 3.13 \\
benzylamine & face & 3.0 & 0/5 & -- & -- & 0.61 \\
benzylamine & face & 3.5 & 2/5 & 47.1 & 52.3 & 0.89 \\
benzylamine & face & 4.0 & 2/5 & 5.2 & 7.9 & 1.31 \\
benzylamine & face & 4.5 & 3/5 & 6.1 & 35.2 & 1.77 \\
benzylamine & face & 5.0 & 5/5 & 4.2 & 21.1 & 2.25 \\
benzylamine & face & 5.5 & 5/5 & 1.0 & 8.3 & 2.73 \\
benzylamine & face & 6.0 & 5/5 & 0.9 & 2.0 & 3.22 \\
benzylamine & edge & 3.0 & 0/3 & -- & -- & 1.00 \\
benzylamine & edge & 3.5 & 0/3 & -- & -- & 1.00 \\
benzylamine & edge & 4.0 & 0/3 & -- & -- & 0.92 \\
benzylamine & edge & 4.5 & 0/3 & -- & -- & 0.82 \\
benzylamine & edge & 5.0 & 0/3 & -- & -- & 1.00 \\
benzylamine & edge & 5.5 & 0/3 & -- & -- & 1.36 \\
benzylamine & edge & 6.0 & 1/3 & 0.4 & 0.4 & 1.78 \\
\bottomrule
\end{tabular}

\end{table}

\section{Driving force}
Adiabatic one-electron potentials with the protocol of the main text in water
($\varepsilon = 78.4$, absolute SHE potential 4.44~V\cite{Jonsson1996}), compared
with measured values\cite{Jonsson1996,Mayhew1999} (Table~\ref{tab:pot}), and
the raw DFT driving force in low-dielectric continua with a contact
Coulomb term $-e^2/(4\pi\varepsilon_0\varepsilon r)$ (Table~\ref{tab:ddft}).
No zero-point or thermal corrections are included.

\begin{table}[h]
\caption{One-electron potentials in water.}
\label{tab:pot}
\centering
\begin{tabular}{lrrr}
\toprule
couple & DFT (V vs NHE) & experiment & error (V) \\
\midrule
Me$_2$NH$^{\bullet+}$/Me$_2$NH & +0.73 & +1.300 & -0.57 \\
MeNH$_2^{\bullet+}$/MeNH$_2$ & +1.26 & -- & -- \\
BnNH$_2^{\bullet+}$/BnNH$_2$ & +1.45 & -- & -- \\
Fl/Fl$^{\bullet-}$ (exp.: FMN, pH 7) & -1.16 & -0.313 & -0.84 \\
\bottomrule
\end{tabular}

\end{table}

\begin{table}[h]
\caption{Raw DFT driving force of the contact pair.}
\label{tab:ddft}
\centering
\begin{tabular}{llrrrr}
\toprule
continuum & amine & separated (eV) & $r=4$ \AA & 5 \AA & 6 \AA \\
\midrule
$\varepsilon=4$ & methylamine & 3.55 & 2.65 & 2.83 & 2.95 \\
$\varepsilon=4$ & benzylamine & 3.59 & 2.69 & 2.87 & 2.99 \\
$\varepsilon=10$ & methylamine & 2.83 & 2.47 & 2.55 & 2.59 \\
$\varepsilon=10$ & benzylamine & 2.97 & 2.61 & 2.69 & 2.73 \\
\bottomrule
\end{tabular}

\end{table}

\section{Cryptochrome}
Orientation-averaged, perpendicular and parallel response of the cryptochrome
pair over the exchange prefactor $J_0$ (both signs of $J_{\rm EH}$;
$J = -2J_{\rm EH}$, $J_{\rm EH} = J_0e^{-\beta r}$, $\beta = 14$~nm$^{-1}$,
$r = 1.90$~nm\cite{Efimova2008}), $k_S = k_P = 1~\mu$s$^{-1}$, coherent
(Table~\ref{tab:cry}).

\begin{table}[h]
\caption{Cryptochrome response (\%).}
\label{tab:cry}
\centering
\small\begin{tabular}{rrrrr}
\toprule
$J_0$ ($\mu$T) & $J$ (MHz) & powder (\%) & $\perp$ (\%) & $\parallel$ (\%) \\
\midrule
8.00\text{e}12 & -1.26 & +0.417 & +0.494 & +0.027 \\
8.00\text{e}12 & +1.26 & +0.386 & +0.464 & -0.023 \\
1.01\text{e}13 & -1.58 & +0.395 & +0.469 & +0.015 \\
1.01\text{e}13 & +1.58 & +0.384 & +0.457 & -0.030 \\
1.27\text{e}13 & -1.99 & +0.371 & +0.454 & -0.005 \\
1.27\text{e}13 & +1.99 & +0.392 & +0.461 & -0.034 \\
1.60\text{e}13 & -2.51 & +0.403 & +0.498 & +0.005 \\
1.60\text{e}13 & +2.51 & +0.413 & +0.491 & -0.029 \\
2.01\text{e}13 & -3.16 & +0.470 & +0.561 & +0.021 \\
2.01\text{e}13 & +3.16 & +0.435 & +0.536 & -0.020 \\
2.53\text{e}13 & -3.98 & +0.573 & +0.681 & +0.027 \\
2.53\text{e}13 & +3.98 & +0.444 & +0.580 & -0.012 \\
3.18\text{e}13 & -5.01 & +0.634 & +0.751 & +0.042 \\
3.18\text{e}13 & +5.01 & +0.404 & +0.483 & +0.007 \\
4.01\text{e}13 & -6.30 & +0.559 & +0.622 & +0.034 \\
4.01\text{e}13 & +6.30 & +0.334 & +0.407 & +0.025 \\
5.05\text{e}13 & -7.93 & +0.383 & +0.420 & +0.061 \\
5.05\text{e}13 & +7.93 & +0.242 & +0.320 & -0.005 \\
6.35\text{e}13 & -9.99 & +0.441 & +0.534 & +0.033 \\
6.35\text{e}13 & +9.99 & +0.242 & +0.299 & +0.049 \\
8.00\text{e}13 & -12.57 & +0.492 & +0.589 & +0.015 \\
8.00\text{e}13 & +12.57 & +0.168 & +0.234 & +0.021 \\
1.01\text{e}14 & -15.83 & +0.282 & +0.372 & +0.024 \\
1.01\text{e}14 & +15.83 & +0.153 & +0.174 & +0.081 \\
1.27\text{e}14 & -19.93 & +0.085 & +0.115 & -0.060 \\
1.27\text{e}14 & +19.93 & +0.150 & +0.131 & +0.126 \\
1.60\text{e}14 & -25.09 & +0.061 & +0.078 & +0.020 \\
1.60\text{e}14 & +25.09 & +0.185 & +0.119 & +0.214 \\
2.01\text{e}14 & -31.58 & +0.044 & +0.049 & +0.023 \\
2.01\text{e}14 & +31.58 & +0.055 & +0.078 & +0.063 \\
2.53\text{e}14 & -39.76 & -0.009 & -0.004 & -0.013 \\
2.53\text{e}14 & +39.76 & +0.036 & +0.041 & +0.019 \\
3.18\text{e}14 & -50.05 & +0.015 & +0.017 & +0.009 \\
3.18\text{e}14 & +50.05 & +0.068 & +0.063 & +0.041 \\
4.01\text{e}14 & -63.01 & +0.010 & +0.011 & +0.007 \\
4.01\text{e}14 & +63.01 & +0.052 & +0.045 & +0.031 \\
5.05\text{e}14 & -79.33 & +0.004 & +0.005 & +0.003 \\
5.05\text{e}14 & +79.33 & +0.039 & +0.049 & +0.015 \\
6.35\text{e}14 & -99.87 & +0.003 & +0.003 & +0.002 \\
6.35\text{e}14 & +99.87 & +0.034 & +0.036 & +0.019 \\
8.00\text{e}14 & -125.73 & +0.001 & +0.002 & +0.001 \\
8.00\text{e}14 & +125.73 & +0.011 & +0.017 & +0.002 \\
\bottomrule
\end{tabular}

\end{table}

\bibliographystyle{unsrt}
\bibliography{refs}

\end{document}
